\documentclass[12pt]{article}

\usepackage[utf8]{inputenc}
\usepackage[left=2.5cm, top=2cm, right=2cm, bottom=2.5cm]{geometry}
\usepackage[most]{tcolorbox}
\usepackage{graphicx}
\usepackage{wasysym}
\usepackage{latexsym}
\usepackage{algorithm}
\usepackage{algpseudocode}
\usepackage{svg}
\usepackage{amsmath}
\usepackage{minted}
\usepackage{caption}
\usepackage{url}
\usepackage[colorlinks]{hyperref}
\hypersetup{
  urlcolor=blue
}

\title{Database Administration \\ Final Project}
\author{Andrés Oswaldo Calderón Romero, Ph.D.}
\date{\today}

\begin{document}

\maketitle

\section{Introduction}

The rapid transformation of energy infrastructure in Colombia calls for innovative approaches to identify suitable locations for the deployment of renewable energy solutions. As part of a broader national strategy to promote sustainable development and improve energy access in underserved regions, the Research and Information Office of the Unidad de Planeación Minero Energética (UPME) has launched a strategic initiative focused on evaluating the feasibility of solar energy in selected territories.

This project supports UPME’s objective by designing and implementing a reproducible geospatial analysis workflow to estimate the solar energy potential of building rooftops across prioritized municipalities. In particular, the project emphasizes the analysis of PDET (Programas de Desarrollo con Enfoque Territorial) territories, which represent a key focus area for post-conflict development and infrastructure enhancement in Colombia.

Leveraging openly available geospatial datasets, the project aims to estimate \textbf{the number and area of building rooftops suitable for solar panel installation}. These datasets, encompassing billions of building outlines derived from high-resolution satellite imagery, provide a valuable resource for quantifying potential energy-harvesting surfaces in urban and rural contexts.

To align with the UPME’s technological infrastructure and modernization goals, the methodology is designed using NoSQL data solutions, enabling scalable storage, efficient querying, and flexible spatial operations. The outcome of this work will be \textbf{a technical report outlining the methodology, results, and recommendations} for identifying optimal locations for proof-of-concept solar farms in the selected regions.

By integrating modern data science tools with real-world energy policy needs, this project bridges technical innovation and strategic planning in support of Colombia’s energy transition and territorial equity.

\section{Problem Statement and Scope}

The Research and Information Office of UPME in Colombia is seeking \textit{a strategic partner} to support the initial phases of a major initiative aimed at identifying potential locations for proof-of-concept projects involving alternative energy solutions. Among these, UPME plans to assess the feasibility of solar farms in various regions across the country.

In order to select potential locations, stakeholders aim to prioritize local governments (municipalities) that have been designated as PDET\footnote{\textit{¿Qué son los PDET?}. 2026. Agencia de Renovación del Territorio. Available at \url{https://www.renovacionterritorio.gov.co/central-pdet/que-son-los-pdet}} territories. The primary goal of the project is to determine the number of buildings within each municipality and to estimate the total rooftop area suitable for solar panel installation. The greater the rooftop area in a municipality, the higher the potential for collecting clean energy through solar panels installed on those roofs. UPME requires a \textbf{detailed report} outlining a reproducible methodology for counting the number of rooftops and aggregating their total area for each PDET municipality evaluating different datasets.

Given the new technologies and infrastructure acquired by the Research and Information Office, a key requirement is the application of NoSQL solutions in the proposed methodology. Additionally, the Office mandates the evaluation of open datasets to query and identify buildings within the target territories, with the goal of comparing the outputs from each source. Three open datasets have been selected as the primary sources for this study, you should pick at least two of them:

\begin{itemize}
  \item \textbf{Microsoft Building Footprints:} Bing Maps has released open building footprint data covering various parts of the world. The dataset includes over 999 million building detections derived from Bing Maps imagery collected between 2014 and 2021, incorporating sources such as Maxar and Airbus. The data is freely available under the Open Data Commons Open Database License (ODbL). More information is available at: \url{https://planetarycomputer.microsoft.com/dataset/ms-buildings}.

  \item \textbf{Google Open Buildings:} This dataset contains 1.8 billion building detections spanning an inference area of 58 million~km\textsuperscript{2}, covering Africa, South Asia, Southeast Asia, Latin America, and the Caribbean. Currently in its third version, the dataset is distributed under the Creative Commons Attribution (CC BY-4.0) and Open Data Commons Open Database License (ODbL) v1.0. More information is available at: \url{https://sites.research.google/gr/open-buildings/}.

  \item \textbf{GlobalBuildingAtlas:} Developed by the Technical University of Munich (TUM), this dataset comprises 2.75 billion building models, representing all structures captured in satellite imagery during 2019. As the most comprehensive collection of its kind, it provides 3D models with a spatial resolution of 3$\times$3 meters. The dataset is distributed under the MIT License with a Commons Clause Restriction. Further information is available at: \url{https://github.com/zhu-xlab/GlobalBuildingAtlas}.

\end{itemize}

In addition to the aforementioned datasets, the Colombian National Administrative Department of Statistics (DANE) provides the \textit{Marco Geoestadístico Nacional} (MGN)\footnote{MGN Manueal de usuario v. 2.0. 2020. DANE.  Available at \url{https://geoportal.dane.gov.co/descargas/descarga_mgn/Manual_MGN.pdf}}, which offers administrative boundaries at various levels, including all Colombian municipalities\footnote{Descarga de datos geoestadísticos. 2025. DANE. Available at \url{https://geoportal.dane.gov.co/servicios/descarga-y-metadatos/datos-geoestadisticos/?cod=111}. Note: At the time of writing, the latest available full version is ``\textit{Versión MGN2025-Colombia. Todos los niveles geográficos, 2025, ZIP, 1.5 GB}''}. Consistent with the scope of this study, the municipal layer should be extracted and subsequently filtered to include only those jurisdictions designated as PDET territories.

\section{Requested Deliverables}
Deliverables are to be submitted weekly and defended during sessions in Weeks 2 (follow-up defense -- 1 random person), and 4 (final defense -- 2 random people). All submissions must be made via commits to each group's GitHub repository and PDF submission via BrightSpace\texttrademark. It is expected that \textbf{ALL} team members contribute equally to every submission; consequently, submissions where the majority of the content is pushed by few members and/or at the last minute will not be considered.

\subsection{Schedule}
  \textit{\textbf{Note:} Please refer to the planner in the \texttt{resources} folder.}
  \begin{enumerate}
    \item \textbf{NoSQL Database Schema Design and Implementation Plan}: This deliverable focuses on meeting the key requirement of using NoSQL solutions for scalable storage and efficient spatial operations. Expected delivery is on \textbf{Week 1} before 2:00 PM in the GitHub repository and BS\texttrademark. It will establish the foundational data structure. This deliverable should include:
    \begin{itemize}
      \item Implementation Plan: A concise but complete description of the methodology to follow.
      \item Schema Design \& Data Modeling: An E-R diagram and a brief data dictionary for the key entities and the interactions among them.
    \end{itemize}

    \item \textbf{PDET Municipality Boundaries Dataset Integration}: This deliverable ensures the project focuses only on the mandated PDET territories, using the required DANE/MGN administrative boundaries. Expected delivery is on \textbf{Week 2} before 2:00 PM in the GitHub repository and BS\texttrademark. This will be evaluated based on:
        \begin{itemize}
          \item Data Acquisition \& Verification: Documentation of download mechanisms and data quality checks.
          \item Data Integrity \& Format: Ensuring the data meets the requirements for ingestion into the database.
          \item NoSQL Spatial Integration: Evidence of implemented spatial indices.
        \end{itemize}

    \item \textbf{Building Footprint Data Loading and Integration}: This task addresses the mandate to evaluate and compare outputs from the Microsoft, Google, or TUM open datasets. Expected delivery is on \textbf{Week 3} before 2:00 PM in the GitHub repository and BS\texttrademark. This deliverable should include:
        \begin{itemize}
          \item The same requirements as the previous deliverable (including spatial indexing) applied to the two selected datasets.
          \item Data Loading Efficiency: An explanation of the techniques used to handle large datasets.
          \item Brief Exploratory Data Analysis (EDA): Highlights of the main statistics and characteristics of the available data.
        \end{itemize}

    \item \textbf{Reproducible Geospatial Analysis Workflow (Rooftop Count and Area Estimation)}: This is the core task, fulfilling the primary goal of the project to estimate the number and total area of building rooftops in each PDET municipality. Expected delivery is on \textbf{Week 4} before 2:00 PM in the GitHub repository and BS\texttrademark. This deliverable should include:
        \begin{itemize}
          \item Reproducibility \& Methodology: A clear explanation of how the required metrics were computed.
          \item Output Data Structure (Tables and Maps): Presentation of the results using appropriate visualizations.
        \end{itemize}

    \item \textbf{Final Technical Report}: This comprehensive and self-contained report is the final deliverable, summarizing the entire project and synthesizing the required data for UPME. Expected delivery is on \textbf{Week 4} before 2:00 PM in the GitHub repository and BS\texttrademark.  A private defense in class will hava place on Monday that week. Only \textbf{two} members of each team will be randomly selected to defend the work. This deliverable will be evaluated based on:
        \begin{itemize}
          \item Documentation of the entire process.
          \item Presentation of results and data visualizations.
          \item Discussion and recommendations.
          \item Overall clarity and completeness of the report.
        \end{itemize}
\end{enumerate}

\section{Conclusions}
This project presents the framework for an analysis aimed at addressing a real-world challenge: the identification and quantification of solar energy potential in developing regions of Colombia. By enabling the loading, processing, and analysis of geospatial data, the proposed solution facilitates a comprehensive evaluation of potential locations through a NoSQL-based approach. The inclusion of both theoretical and empirical analyses ensures a robust assessment of the available data and contextual conditions. The final deliverable, \textbf{a detailed technical report}, will serve as a solid foundation for selecting the most appropriate methodologies and locations in future implementation phases.

\end{document}
